\section{Problem Description}
\label{sec:problem}

\subsection{System Model}

We consider a cloud service operator managing a portfolio of \emph{delay-tolerant} AI workloads distributed across $|\mathcal{R}| = 5$ geographically distributed datacenter regions.
Each region $r \in \mathcal{R}$ is connected to a distinct electricity grid with its own carbon intensity profile.
The planning horizon spans $|\mathcal{T}| = 17{,}544$ hourly time steps covering calendar years 2024 and 2025.

\textbf{Workload model.}
Total compute demand $D$ (in kWh) must be satisfied over the planning horizon.
At each hour $t$, each region $r$ can accept up to $C_{\max,r}$ kWh of compute load.
A fraction of the load is \emph{flexible}: it can be shifted in time (within a delay budget) or routed to a different region.
Inflexible (latency-sensitive) load is excluded from the scheduling decision.

\textbf{Grid regions.}
The five regions are selected to represent structurally distinct grid types:

\begin{table}[h]
\centering
\begin{tabular}{llp{5.5cm}l}
\toprule
\textbf{Region} & \textbf{Zone ID} & \textbf{Grid character} & \textbf{CI profile} \\
\midrule
PJM       & US-MIDA-PJM & Coal/gas/nuclear mix; high demand variability & High, volatile \\
NYISO     & US-NY-NYIS  & Hydro + nuclear dominant; some gas peaking & Medium, diurnal \\
Finland   & FI          & Hydro + nuclear + wind; Nordic interconnect  & Low, seasonal \\
Belgium   & BE          & High nuclear share; European grid coupling  & Medium-low \\
Singapore & SG          & Near-total natural gas; isolated island grid & High, stable \\
\bottomrule
\end{tabular}
\caption{Target grid regions and their structural characteristics.}
\label{tab:regions}
\end{table}

Singapore serves as a critical control case: its low CI variability (CV $\approx$ \todo{fill after EDA}) means temporal shifting yields minimal benefit, validating that the model correctly assigns near-zero temporal gains to low-variability grids.

\subsection{Data}

\textbf{Carbon intensity (CI).}
Hourly CI values (gCO\textsubscript{2}eq/kWh) are retrieved from the ElectricityMaps API via the \texttt{/v3/carbon-intensity/past} endpoint.
ElectricityMaps aggregates raw grid measurements from Transmission System Operators (TSOs) and applies machine-learning models to fill short gaps, producing a continuous, pre-validated hourly series~\cite{radovanovic2023}.
The dataset covers 2024-01-01 00:00 UTC to 2025-12-31 23:00 UTC, yielding 17,544 observations per region with no missing values.

\textbf{Renewable fraction (RF).}
Hourly RF values are computed from the \texttt{/v3/power-breakdown/past} endpoint as:
\[
  \mathrm{RF}_r(t) = \frac{\sum_{s \in \mathcal{S}_{\mathrm{clean}}} P_{r,s}(t)}{\sum_{s} P_{r,s}(t)},
  \quad \mathcal{S}_{\mathrm{clean}} = \{\text{wind, solar, hydro, geothermal, biomass}\}
\]
Nuclear is excluded from the renewable set: while low-carbon, it is dispatchable and not intermittent, making it a distinct scheduling consideration.

\textbf{Electricity price (budget parameter).}
Per-region electricity price series are used to parameterize the budget constraint C4 (Section~\ref{sec:formulation}).
Sources: PJM Locational Marginal Prices (LMPs), NYISO day-ahead prices, ENTSO-E (Finland and Belgium), and EMC spot prices (Singapore).
\todo{Confirm data availability and add source URLs.}

\subsection{Key Indicators}

Three indicators drive the model and analysis:

\begin{itemize}[leftmargin=1.5em]
  \item \textbf{CI$_r(t)$}: the carbon cost of one kWh of compute in region $r$ at hour $t$.
        This is the primary optimization signal.
  \item \textbf{RF$_r(t)$}: the renewable fraction of power consumption at time $t$.
        Incorporated in the objective to incentivize absorption of variable renewables beyond simply targeting low-CI periods.
  \item \textbf{CV$_r$}: the coefficient of variation of CI for region $r$, computed as $\sigma(\mathrm{CI}_r) / \mu(\mathrm{CI}_r)$ over the full horizon.
        Used as a summary diagnostic of each region's scheduling opportunity.
\end{itemize}

\subsection{Assumptions}

\begin{enumerate}[leftmargin=1.5em]
  \item \textbf{Perfect historical information.}
        The LP is solved with full knowledge of the historical CI and RF series (backtesting mode).
        This yields an upper bound on achievable carbon savings; real-time deployment would require CI forecasts.
  \item \textbf{Separable compute-to-power conversion.}
        One unit of compute load draws one unit of power, abstracting away PUE (Power Usage Effectiveness) differences across regions.
        Regional PUE heterogeneity is a direction for future work.
  \item \textbf{No cross-region data transfer cost.}
        Network latency and data movement costs are not modeled; the model is appropriate for batch workloads where data locality is not binding.
  \item \textbf{Stationary demand $D$.}
        Total demand is fixed and must be satisfied in full.
        Load shedding or demand reduction are not considered.
\end{enumerate}
